\clearpage
\section{Evaluating and sampling language models}\label{sec:evaluation}
\subsection*{One total loss, several units}
Let $S=-\sum\log p_\theta(x_t\mid x_{<t})$ be total negative log likelihood
in nats, $N$ the number of scored predictions, and $B_{\mathrm{utf8}}$ the
number of evaluated text bytes. The same score can be reported as
\[
 \ell=\frac SN,\quad \mathrm{BPT}=\frac{S}{N\log2},\quad
 \mathrm{PPL}=e^{S/N},\quad\mathrm{BPB}=\frac{S}{B_{\mathrm{utf8}}\log2}.
\]
Perplexity is the inverse geometric mean of the probabilities assigned to
observed targets. For a uniform distribution on ten digits, every target
has probability $1/10$, so every nonempty sequence has perplexity ten.
For a general model it is an effective branching factor, not an actual
count of equally likely options at every position.

\fig{metric-units}{Start from a common numerator. The denominator determines
whether a quantity is per token or per byte; exponentiation produces perplexity.}

\fig{tokenizer-metrics}{A constructed comparison with four total bits of loss
on the same four-byte text. Fewer scored tokens increase per-token perplexity
here; bits per byte stays one.}

The figure holds total loss fixed to isolate units. Real tokenizers also
change the model's events and fitted probabilities, so BPB can change too.
Cross-tokenizer BPB comparisons require identical evaluated text,
normalization, document boundaries, and special-token scoring conventions.

\textbf{End tokens affect the numerator and token count.} If a one-byte
document \texttt{a} has probability $1/2$ for its character and $1/2$ for
\EOS{}, then $N=2$, $S=2\log2$, and $B_{\mathrm{utf8}}=1$. Perplexity
is two and BPB is two. The boundary token contributes loss without adding
a literal byte to the evaluated text.

Use training loss to diagnose fitting, development loss to make choices,
and a fixed test set for final comparisons. Task-based evaluation asks a
different question: a model can improve average token loss while becoming
worse on a particular task or subgroup.

\selfcheck{Why does averaging the perplexities of two batches generally differ
from exponentiating their combined mean loss? Accumulate total loss and
prediction counts first, especially when batches have different sizes.}

\nextpage{Sampling, data filters, and experimental conclusions}
To sample a categorical distribution, draw $u$ uniformly from $[0,1)$ and
find the interval containing it in the cumulative probabilities. For a
language model, sample a token, append it to the context, recompute the
distribution, and repeat until \EOS{} or a stated length cap.

\fig{sampling}{A categorical draw for $(1/2,1/3,1/6)$. A draw at $u=0.7$
selects the second interval; generation repeats the same operation with updated context.}

For an interpolated model, either sample from the combined distribution or
first draw a component according to its mixture weight and then sample
from that component. Both implement the same probability law. Sampling
from raw counts while reporting the likelihood of a smoothed mixture
would describe two different models. Temperature and top-$k$ truncation
also change the generation distribution and should be labeled.

\fig{pipeline-experiments}{Two model-based data workflows in L02. A filter makes
a selection from scored documents; a sample/retrain loop changes the data
used by the next model. Neither diagram establishes a benefit on its own.}

\textbf{A filter score is a signal.} A reference LM may give lower loss to
text resembling its own training distribution. Selecting low-loss documents
can therefore change topic, language, and style. The course's reference
trigram is a teaching analogue of perplexity-based filtering; a threshold
does not certify truth, usefulness, or broad data quality.

\textbf{Read the budget with the curve.} In the saved L02 sample/retrain
demonstration, held-out BPB rises from 1.299 to 1.940 over five rounds.
The original training corpus has 6,076,505 tokens; the final synthetic
round has 100,681. Each round samples 2,000 documents capped at 128 tokens,
so length caps and changing token budgets are confounded with replacing
real data by samples. The observation is deterioration in this run, not
an isolated causal estimate of one mechanism.

An informative follow-up would hold total training tokens, vocabulary,
evaluation text, and training procedure fixed, then compare real-only,
synthetic-only, and mixed data across several seeds. This is an experiment
design, not a result reported by the notes.

\takeaway{Scoring, sampling, and selecting training data are distinct
operations. State the distribution and the comparison each one uses.}
\selfcheck{If a sampler stops at a length cap without \EOS{}, should the
result be recorded as naturally terminated? Keep the two stopping reasons distinct.}
